\section{Actual analytic superlevels of the canonical boundary count}

The maximal-order step of W{\l}odarczyk\cite{wlodarczyk05-hironaka},
proof of Proposition \texttt{pr:resol}, Step 1a, uses the number of
original boundary members through a point. Its superlevel is globally
analytic for a genuine analytic boundary. This supplies another actual
canonical invariant layer needed before the analytic extension construction
of Bierstone--Milman\cite{bierstone-milman97-invent-math}, Section 13.

\begin{lemma}[Actual finite intersection dictionary]
\label{hir:boundary-count-superlevel-intersection-dictionary}
\lean{Hironaka.HironakaBoundaryCountAnalyticSuperlevels.boundaryCount_superlevel_eq_biUnion}
\leanok
For a finite list $E$ and $s\in\mathbb N$,
\[
 \{x:s\le \#\{i:x\in E_i\}\}
 =\bigcup_{S\subset\{1,\ldots,|E|\},\ |S|=s}\ \bigcap_{i\in S} E_i.
\]
Indices retain the actual original list occurrences. No choice of
intersections or superlevel certificate is part of the input.
\end{lemma}
\begin{proof}
At a point of count at least $s$, the finite set of indices through that
point has a subset $S$ of cardinality exactly $s$. The point lies in every
member indexed by $S$, so it lies in the displayed union. Conversely,
membership in a displayed intersection puts all $s$ indices of $S$ into
the actual index set through the point; cardinality monotonicity proves
the count bound. For $s=0$, the empty intersection is the whole ambient
space. For $s>|E|$, both sides are empty. Repeated list members do not
invalidate this indexed combinatorial identity.
\end{proof}

\begin{theorem}[Produced analytic boundary-count superlevels]
\label{hir:boundary-count-global-analytic-superlevels}
\uses{hir:boundary-count-superlevel-intersection-dictionary}
\lean{Hironaka.HironakaBoundaryCountAnalyticSuperlevels.isComplexAnalyticSubset_boundaryCount_superlevel}
\leanok
If every actual boundary member is analytic, each boundary-count
superlevel is globally analytic in every ambient dimension, without
compactness.
\end{theorem}
\begin{proof}
Finite intersections of the genuine analytic members are analytic,
including the empty intersection, the whole ambient space. The exact
dictionary expresses the superlevel as a finite union of those actual
intersections. Finite unions are analytic by products of their actual
local defining equations. Thus the analytic structure is produced from
the original members, rather than inferred from closedness or topological
upper semicontinuity.
\end{proof}

\begin{theorem}[Actual transformed original-member level superlevels]
\label{hir:actual-step-one-level-analytic-superlevels}
\uses{hir:boundary-count-global-analytic-superlevels,hir:step-one-seq}
\lean{Hironaka.HironakaBoundaryCountAnalyticSuperlevels.isComplexAnalyticSubset_stepOne_level_superlevel}
\leanok
Let $\mathfrak J$ be a genuinely well formed marked ideal in positive
ambient dimension, and $T$ a genuine admissible test tower of its actual
homogenized coefficient marked ideal $\overline{\mathfrak J}$. The locus
where the genuine original-member level $s_T$ is at least $s$ is globally
analytic.
\end{theorem}
\begin{proof}
The actual homogenized coefficient marked ideal is well formed. Its
genuine transform along $T$ therefore has an actual SNC boundary. Taking
the original number of its members gives the genuine list of original
strict transforms; this sublist remains SNC. In positive dimension,
every actual SNC member is a closed smooth hypersurface, as follows by
the genuine singleton member-intersection regularity theorem. Hence each
is analytic. The definition of $s_T$ is exactly the actual boundary count
of this list, so apply the produced analytic-superlevel theorem. No
analytic presentation of the level or transformed members is requested
from the caller.
\end{proof}

\begin{corollary}[Actual masked StepOne first-entry analytic superlevels]
\label{hir:actual-masked-step-one-first-entry-analytic-superlevels}
\uses{hir:actual-step-one-level-analytic-superlevels,hir:step-one-seq}
\lean{Hironaka.HironakaBoundaryCountAnalyticSuperlevels.isComplexAnalyticSubset_stepOneInvM_first_superlevel}
\leanok
Under the same genuine geometric inputs, the first-entry superlevels of
the actual masked StepOne invariant are globally analytic.
\end{corollary}
\begin{proof}
At special points the actual invariant is $(s_T,\infty)$; at every other
point it is $(s_T,\text{actual lower value})$. Consequently its first
entry equals $s_T$ at every point. Substitute this identity into the
preceding theorem. This connection uses the actual invariant constructor,
not a supplied tuple or an analyticity certificate. The lower/history
entries and the full auxiliary-chain invariant still require separate
analyticity and smoothification producers before general noncompact
global resolution can be claimed.
\end{proof}
